\documentclass[11pt]{article}

\usepackage[margin=1in]{geometry}
\usepackage{amsmath,amssymb,amsthm}
\usepackage{bm}
\usepackage{hyperref}
\usepackage{parskip}

\renewcommand{\ne}{n_{e}}  % shadow LaTeX kernel's \ne ("not equal") -- not used in this doc
\newcommand{\nfc}{n_{\mathrm{FC}}}
\newcommand{\ncx}{n_{\mathrm{CX}}}
\newcommand{\Vfc}{\lvert V_{\mathrm{FC}}\rvert}
\newcommand{\Vcx}{\lvert V_{\mathrm{CX}}\rvert}
\newcommand{\Vfcr}{V_{\mathrm{FC},r}}
\newcommand{\Vcxr}{V_{\mathrm{CX},r}}
\newcommand{\Si}{S_{i}}
\newcommand{\Scx}{S_{\mathrm{CX}}}
\newcommand{\Sreact}{S}
\newcommand{\Dped}{D_{\mathrm{ped}}}
\newcommand{\fFC}{f_{\mathrm{FC}}}
\newcommand{\fCX}{f_{\mathrm{CX}}}
\newcommand{\fsa}[1]{\left\langle #1 \right\rangle}
\newcommand{\gradrsq}{\fsa{|\nabla r|^{2}}}
\newcommand{\neginf}{x=-\infty}
\newcommand{\dneneg}{\left.\dfrac{d\ne}{dx}\right|_{\neginf}}
\newcommand{\dnexnow}{\dfrac{d\ne}{dx}}
\newcommand{\xin}{x_{\mathrm{inner}}}
\newcommand{\nFCzero}{\fsa{\nfc}(0)}
\newcommand{\Afactor}{\Vfc\,\fFC}
\newcommand{\Bfactor}{\Vcx\,\fCX}
\newcommand{\Qfac}{\mathcal{Q}}
\newcommand{\Ccx}{C_{\mathrm{CX}}}
\newcommand{\Ecx}{E_{\mathrm{CX}}}

\title{Derivation of Equation (15) in Saarelma \emph{et al.}~2023\\
  \large Nucl.\ Fusion \textbf{63} 052002 --- full form including CX neutrals}
\author{}
\date{}

\begin{document}
\maketitle

This document derives Equation (15) of the Saarelma--Connor pedestal
neutral-penetration paper.  It is the full form \emph{with} charge-exchange
neutrals retained, of which Eq.~(16) is the $\fsa{\ncx}\to 0$ limit.
All the preparatory machinery (three-fluid model, flux-surface averaging,
form factors, Eqs.~(8)--(10)) is identical to \texttt{derivation\_eq16.tex}
and is summarized here.  The additional steps unique to Eq.~(15) are
(i)~explicit integration of~\eqref{eq:8} to obtain $\fsa{\nfc}(x)$ as
an exponential, and (ii)~algebraic elimination of $\fsa{\ncx}$ via the
first-integral~\eqref{eq:10}.

\section{Starting point: the three Eqs.~(8), (9), (10)}
\label{sec:start}

From \texttt{derivation\_eq16.tex} (reproduced here for reference):
\begin{align}
  \ne\Si\fsa{\nfc}
    &= \frac{\Afactor}{1+\Scx/\Si}\,\frac{d\fsa{\nfc}}{dx}
     = \frac{\Afactor\,\Si}{\Si+\Scx}\,\frac{d\fsa{\nfc}}{dx}, \label{eq:8}\\[4pt]
  \ne\Si\fsa{\ncx}
    &= \Bfactor\,\frac{d\fsa{\ncx}}{dx}
     + \frac{\Afactor\,\Scx}{2(\Si+\Scx)}\,\frac{d\fsa{\nfc}}{dx}, \label{eq:9}\\[4pt]
  \gradrsq\,\Dped\,\bigl(\dnexnow - \dneneg\bigr)
    &= -\,\frac{\Si+\Scx/2}{\Si+\Scx}\,\Afactor\,\fsa{\nfc}
       \;-\; \Bfactor\,\fsa{\ncx}. \label{eq:10}
\end{align}
The LHS of~\eqref{eq:10} is the transport flux relative to its inner-edge
value: at $x\to-\infty$ both neutral densities vanish and
$\gradrsq\Dped\dnexnow\to\gradrsq\Dped\dneneg$, which fixes the integration
constant.  Assumptions A1--A3 (locally mono-energetic neutrals, slowly
varying geometry and form factors, slab-like neutral closure) have
already been used to arrive here; see the Eq.~(16) document for the full
list.

The electron continuity equation in its first non-integrated form is
\begin{equation}
  \frac{d}{dx}\!\left[\gradrsq\,\Dped\,\dnexnow\right]
    \;=\; -\,\ne\,\Si\,\bigl(\fsa{\nfc}+\fsa{\ncx}\bigr),
  \label{eq:4-1D}
\end{equation}
i.e.\ Eq.~(4) of the paper after dividing by $\oint R^{2}d\theta$.  Our
task is to use~\eqref{eq:8}--\eqref{eq:10} to rewrite the RHS
of~\eqref{eq:4-1D} purely in terms of $\ne(x)$, $\dnexnow$, and known
quantities.  That is exactly what Eq.~(15) of the paper is.

\section{Step 1: integrate Eq.~(8) for $\fsa{\nfc}$}
\label{sec:integrate-fc}

Rearrange~\eqref{eq:8}:
\begin{equation}
  \frac{d\fsa{\nfc}}{dx} \;=\; \frac{\ne\,(\Si+\Scx)}{\Afactor}\,\fsa{\nfc}.
\end{equation}
This is separable.  Writing $\Afactor=\Vfc\fFC$ explicitly,
\begin{equation}
  \frac{d\ln\fsa{\nfc}}{dx} \;=\; \frac{\ne(x)\,(\Si+\Scx)}{\Vfc\,\fFC}.
\end{equation}
Integrate from $x=0$ (the separatrix) down to an arbitrary $x<0$:
\begin{equation}
  \boxed{\;
    \fsa{\nfc}(x) \;=\; \nFCzero \,
      \exp\!\left[\,\int_{0}^{x}\!\frac{\ne(x')\,(\Si+\Scx)}{\fFC\,\Vfc}\,dx'\,\right].
  \;}
  \label{eq:11}
\end{equation}
The upper limit of the integral is the \emph{current} $x$; for $x<0$
the integral runs ``backwards'' so $\int_{0}^{x}=-\int_{x}^{0}<0$.  The
integrand is positive, so the exponent is negative and
$\fsa{\nfc}(x)<\nFCzero$ for $x<0$ --- FC density decays monotonically
inward, as physics demands.

Equation~\eqref{eq:11} is the paper's Eq.~(11).  It is \emph{closed-form}
given $\ne(x)$: the FC neutral density at any $x$ is determined by
integrating the local electron density against $(\Si+\Scx)/(\fFC\Vfc)$
from the separatrix inward.  This is the ``optical-depth'' picture: the
exponent is the integrated neutral-absorption rate divided by the radial
ballistic speed.

\paragraph{Boundary condition at $x=0$.}  The paper takes
$\nFCzero$ as an input, either from neutral-gas puff modelling or from
fitting the measured pedestal top/bottom densities.  In the code,
\texttt{self.nFC\_x0} is computed in the \texttt{neutrals} module from
the upstream recycling flux.

\paragraph{Assumption A2b revisited.}  When we wrote~\eqref{eq:8} with
$\Vfc\fFC$ outside the $d/dx$ we used the fact that $\Vfc$ is constant
and $\fFC$ varies slowly.  No new assumption here.

\section{Step 2: solve Eq.~(10) for $\fsa{\ncx}$}
\label{sec:solve-ncx}

This is purely algebraic.  From~\eqref{eq:10}:
\begin{equation}
  \Bfactor\,\fsa{\ncx} \;=\; -\,\gradrsq\,\Dped\,\bigl(\dnexnow - \dneneg\bigr)
     \;-\; \frac{\Si+\Scx/2}{\Si+\Scx}\,\Afactor\,\fsa{\nfc}.
\end{equation}
Divide through by $\Bfactor=\Vcx\fCX$:
\begin{equation}
  \boxed{\;
    \fsa{\ncx} \;=\;
      -\,\frac{\gradrsq\,\Dped}{\Bfactor}\,\bigl(\dnexnow - \dneneg\bigr)
      \;-\; \frac{\Afactor}{\Bfactor}\,\frac{\Si+\Scx/2}{\Si+\Scx}\,\fsa{\nfc}.
  \;}
  \label{eq:12}
\end{equation}
This is the closure for the CX neutrals.  Geometrically: the CX density
is driven by two effects,
\begin{itemize}
  \item a flux-divergence term proportional to $(\dnexnow-\dneneg)$
        (the electron pressure-gradient excess over the core),
  \item an FC-born term proportional to $\fsa{\nfc}$ (generation of warm
        neutrals by the CX reaction on FC atoms).
\end{itemize}
Both terms appear with a minus sign on the right, which is correct:
$\fsa{\ncx}\ge 0$ because both $(\dnexnow-\dneneg)$ and $\fsa{\nfc}$ are
non-positive in the pedestal (density gradient is less negative at the
inner boundary; the subtraction gives a non-positive increment) and
positive, respectively.  The signs are discussed further in §\ref{sec:signs}.

\section{Step 3: sum the neutral densities}
\label{sec:sum}

Add~\eqref{eq:11} and~\eqref{eq:12}:
\begin{align}
  \fsa{\nfc} + \fsa{\ncx}
    &= -\,\frac{\gradrsq\,\Dped}{\Bfactor}\bigl(\dnexnow-\dneneg\bigr) \notag\\
    &\quad + \left[\,1 - \frac{\Afactor}{\Bfactor}\cdot\frac{\Si+\Scx/2}{\Si+\Scx}\,\right]\fsa{\nfc}.
  \label{eq:13}
\end{align}
Now substitute~\eqref{eq:11} for $\fsa{\nfc}$ in the bracketed term:
\begin{align}
  \fsa{\nfc} + \fsa{\ncx}
    &= -\,\frac{\gradrsq\,\Dped}{\Bfactor}\bigl(\dnexnow-\dneneg\bigr) \notag\\
    &\quad + \left[\,1 - \frac{\Vfc\fFC(\Si+\Scx/2)}{\Vcx\fCX(\Si+\Scx)}\,\right]
             \nFCzero\,\exp\!\left[\int_{0}^{x}\!\frac{\ne(x')(\Si+\Scx)}{\fFC\,\Vfc}\,dx'\right].
  \label{eq:14}
\end{align}
Equation~\eqref{eq:14} is essentially Saarelma's Eq.~(14).

\section{Step 4: the full Eq.~(15)}
\label{sec:eq-15}

Substitute~\eqref{eq:14} into the electron continuity
equation~\eqref{eq:4-1D}:
\begin{equation}
  \boxed{\;
  \begin{aligned}
    \frac{d}{dx}\!\left[\gradrsq\,\Dped\,\dnexnow\right]
      &= -\,\ne\,\Si\,\Biggl\{
            -\,\frac{\gradrsq\,\Dped}{\Vcx\,\fCX}\bigl(\dnexnow-\dneneg\bigr)  \\
      &\qquad + \left[\,1 - \frac{\Vfc\,\fFC(\Si+\Scx/2)}{\Vcx\,\fCX(\Si+\Scx)}\,\right]
                \nFCzero\,\exp\!\left[\int_{0}^{x}\!\frac{\ne(x')(\Si+\Scx)}{\fFC\,\Vfc}\,dx'\right]
         \Biggr\}.
  \end{aligned}
  \;}
  \label{eq:15-derived}
\end{equation}
This is our fully derived Eq.~(15).  Every approximation used is the one
inherited from deriving Eqs.~(8)--(10).  Comparing with the form in the
paper (and in the attached image of Eq.~(15)),
\begin{equation}
  \frac{d}{dx}\!\left[\gradrsq\,\Dped\,\dnexnow\right]
    = -\,\ne\,\Si\!\left[-\,\Dped\bigl(\dnexnow-\dneneg\bigr)
         + \bigl(1 - \tfrac{\Afactor(\Si+\Scx/2)}{\Bfactor(\Si+\Scx)}\bigr)\,
           \nFCzero\,e^{\int_{0}^{x}\!\ne(\Si+\Scx)/(\fFC\Vfc)\,dx'}\right],
  \label{eq:15-paper}
\end{equation}
we see that the paper's first term is written as $-\Dped\bigl(\dnexnow-\dneneg\bigr)$
whereas our rigorous derivation gives
\[
  -\,\frac{\gradrsq\,\Dped}{\Vcx\,\fCX}\bigl(\dnexnow-\dneneg\bigr).
\]
The difference is a factor of $\gradrsq/(\Vcx\fCX)$ with units of $[s/m]$,
so~\eqref{eq:15-paper} as literally written in the paper is \emph{not
dimensionally consistent}.  The second term has units of
$\mathrm{m^{-3}}$ (as required for $\ne\Si[\cdot]$ to equal a volumetric
rate $\mathrm{m^{-3}/s}$), whereas the paper's first term as written has
units of $\mathrm{m^{-2}/s}$.  The inconsistency is resolved by restoring
the missing $\gradrsq/(\Vcx\fCX)$ factor --- either the paper absorbs it
into its ``$\Dped$'' in that term (shorthand) or this is a typographical
simplification.

\paragraph{Assumption A9 (dimensional-simplification shorthand).}
When writing Eq.~(15) in the compact paper form, we identify
\[
  \frac{\gradrsq\,\Dped}{\Vcx\,\fCX} \;\longrightarrow\; \Dped,
\]
i.e.\ the geometric factor $\gradrsq$ and the CX-neutral kinetic factor
$(\Vcx\fCX)^{-1}$ are treated symbolically, a similar abuse of notation
to what was done in the derivation of~(16).  For code purposes the
\emph{derived} form~\eqref{eq:15-derived} is the physically meaningful
one, and it is the form we use below.

\section{Physical form of the ODE}
\label{sec:phys-ode}

Let $f(x)\equiv\gradrsq\Dped$ (code's \texttt{f\_arr}), and define a
convenient CX-flux coefficient
\begin{equation}
  \Ccx(x) \;\equiv\; 1 \;-\; \frac{\Vfc\,\fFC}{\Vcx\,\fCX}\cdot\frac{\Si+\Scx/2}{\Si+\Scx}.
  \label{eq:Ccx-def}
\end{equation}
This is exactly \texttt{C\_cx} in \texttt{check\_normalization} at
\texttt{solver.py} lines~834--835.  Expand the LHS of~\eqref{eq:15-derived}:
\begin{equation}
  f\,\ne'' + f'\,\ne'
    = -\,\ne\,\Si\!\left[
       -\,\frac{f}{\Vcx\,\fCX}(\ne'-\dneneg) + \Ccx\,\nFCzero\,\mathcal{E}(x)
       \right],
  \label{eq:phys-expand}
\end{equation}
where we have abbreviated the exponential kernel as
\begin{equation}
  \mathcal{E}(x) \;\equiv\; \exp\!\left[\,\int_{0}^{x}\frac{\ne(x')(\Si+\Scx)}{\fFC\,\Vfc}\,dx'\,\right].
  \label{eq:Eker-def}
\end{equation}
Divide through by $f$ and rearrange:
\begin{equation}
  \boxed{\;
    \ne'' \;=\; c_A(x)\,\ne\,\ne'
         \;-\; c_B(x)\,\ne
         \;-\; c_{\Ecx}(x)\,\ne\,\mathcal{E}(x)
         \;-\; c_K(x)\,\ne',
  \;}
  \label{eq:phys-ODE}
\end{equation}
with physical coefficients
\begin{align}
  c_A(x)    &\;=\; \frac{\Si}{\Vcx\,\fCX}, \label{eq:cA}\\
  c_B(x)    &\;=\; \frac{\Si\,\dneneg}{\Vcx\,\fCX}, \label{eq:cB}\\
  c_{\Ecx}(x) &\;=\; \frac{\Si\,\Ccx\,\nFCzero}{f}, \label{eq:cEcx}\\
  c_K(x)    &\;=\; \frac{f'(x)}{f(x)}. \label{eq:cK}
\end{align}

\paragraph{Recovery of Eq.~(16).}
The no-CX limit is $\fsa{\ncx}\to 0$ in~\eqref{eq:10}, which forces
$-\gradrsq\Dped(\dnexnow-\dneneg)=(S_i+S_{CX}/2)/(S_i+S_{CX})\,\Afactor\fsa{\nfc}$
and, after substituting, collapses~\eqref{eq:15-derived} to
Eq.~(16).  The structural difference is that in Eq.~(16) the ``diffusive
flux'' term appears with the $1/\Vfc$ factor (from the FC closure) and
the $(S_i+\Scx)$ enhancement, whereas here in Eq.~(15) it is the CX
fluid that sets the flux relation $\fsa{\ncx}\propto
(\dnexnow-\dneneg)/(\Vcx\fCX)$.

\section{Non-dimensionalization}
\label{sec:nondim}

Introduce the same scales as in the Eq.~(16) document,
\[
  N(\xi)\equiv\ne(x)/n_{0}, \qquad \xi\equiv x/L, \qquad
  n_{0}\equiv\ne(0), \quad L\equiv |\xin|,
\]
so $\ne'=(n_{0}/L)N'$ and $\ne''=(n_{0}/L^{2})N''$.  The Neumann BC is
$(dN/d\xi)_{\xi=-1}=(L/n_{0})\dneneg$ (code's
\texttt{self.dNdxi\_neginf}).

Substitute into~\eqref{eq:phys-ODE} term by term and multiply through by
$L^{2}/n_{0}$:
\begin{equation}
  \boxed{\;
    \frac{d^{2}N}{d\xi^{2}}
      \;=\; A\,N\,\frac{dN}{d\xi}
         \;-\; B\,N
         \;-\; \Ecx\,N\,\mathcal{E}(x(\xi))
         \;-\; K\,\frac{dN}{d\xi},
  \;}
  \label{eq:nondim-ODE}
\end{equation}
with dimensionless coefficients
\begin{align}
  A  &\;=\; n_{0}\,L\,c_A \;=\; \frac{n_{0}\,L\,\Si}{\Vcx\,\fCX}, \label{eq:A}\\[4pt]
  B  &\;=\; L^{2}\,c_B \;=\; \frac{L^{2}\,\Si\,\dneneg}{\Vcx\,\fCX}
          \;=\; \frac{n_{0}\,L\,\Si}{\Vcx\,\fCX}\,
                \left.\frac{dN}{d\xi}\right|_{\xi=-1}, \label{eq:B}\\[4pt]
  \Ecx &\;=\; L^{2}\,c_{\Ecx} \;=\; \frac{L^{2}\,\Si\,\Ccx\,\nFCzero}{f}, \label{eq:Ecx}\\[4pt]
  K  &\;=\; L\,c_K \;=\; \frac{L\,f'(x)}{f(x)} \;=\; \frac{1}{f}\frac{df}{d\xi}. \label{eq:K}
\end{align}
These are the ``rigorous'' (dimensionally-consistent) non-dimensional
coefficients for the iterative solve of Eq.~(15).

\section{Mapping to the code's \texttt{ode\_solv}}
\label{sec:code-map}

The code splits the $N$-linear part of~\eqref{eq:nondim-ODE} into three
pieces (\texttt{K1}, \texttt{K2}, \texttt{K3}), all multiplying $N$, and
groups $B$ and $\Ecx$ differently from the formulation above.
Specifically \texttt{ode\_solv} (lines~960--965) writes
\begin{equation}
  N'' \;=\; A\,N' + B\,N\,N' + (K_{1}+K_{2}+K_{3})\,N,
  \label{eq:code-ODE}
\end{equation}
with
\begin{align}
  A_{\text{code}}        &\;=\; -\,\frac{L\,f'(x)}{f(x)} \;=\; -K, \label{eq:Acode}\\[4pt]
  B_{\text{code}}        &\;=\; \frac{\Si\,\Dped\,n_{0}\,L}{f\,\Vfc}
                              \;=\; \frac{\Si\,n_{0}\,L}{\gradrsq\,\Vfc}, \label{eq:Bcode}\\[4pt]
  K_{1,\text{code}}      &\;=\; -\,\frac{L^{2}\,\Si\,\Dped\,\dneneg}{f\,\Vfc}
                              \;=\; -\,\frac{L^{2}\,\Si\,\dneneg}{\gradrsq\,\Vfc}, \label{eq:K1code}\\[4pt]
  K_{2,\text{code}}      &\;=\; -\,\frac{L^{2}\,\Si\,\nFCzero}{f}\,\mathcal{E}, \label{eq:K2code}\\[4pt]
  K_{3,\text{code}}      &\;=\; +\,\frac{L^{2}\,\Si\,\nFCzero}{f}\,
                                \frac{\Vfc\fFC}{\Vcx\fCX}\,\frac{\Si+\Scx/2}{\Si+\Scx}\,\mathcal{E}. \label{eq:K3code}
\end{align}
Identifying with~\eqref{eq:phys-ODE}:
\begin{itemize}
  \item The code's $A_{\text{code}}$ and the derivation's $-K$ match exactly.
  \item $K_{2,\text{code}}+K_{3,\text{code}}=
        -(L^{2}\Si\,\Ccx\,\nFCzero/f)\,\mathcal{E}=-\Ecx\,\mathcal{E}$ matches the
        third term of~\eqref{eq:nondim-ODE} exactly.
  \item $B_{\text{code}}$ and $K_{1,\text{code}}$ together implement the
        $A\,NN' - B\,N$ piece of~\eqref{eq:nondim-ODE} \emph{but} with
        $1/(\gradrsq\,\Vfc)$ in place of $1/(\Vcx\fCX)$ in $c_A$ and
        $c_B$ (compare \eqref{eq:A}--\eqref{eq:B}).  This corresponds to
        the paper's Eq.~(15) being written with $\Dped$ instead of
        $\gradrsq\Dped/(\Vcx\fCX)$ in the first term --- i.e.\ the code
        implements the paper's \emph{written} form~\eqref{eq:15-paper},
        not the fully-derived form~\eqref{eq:15-derived}.  The mismatch
        factor is $(\Vcx\fCX)/(\gradrsq\Vfc)$.
\end{itemize}

\paragraph{Note on a potential code/paper discrepancy.}  If instead the
code were to use $1/(\Vcx\fCX)$ (the rigorous form), then
\begin{equation}
  B_{\text{code,rigorous}} \;=\; \frac{\Si\,n_{0}\,L}{\Vcx\fCX},
  \qquad
  K_{1,\text{code,rigorous}} \;=\; -\,\frac{L^{2}\,\Si\,\dneneg}{\Vcx\fCX},
\end{equation}
which is dimensionally consistent regardless of the geometry and matches
the Eq.~(10) closure of~\eqref{eq:12} directly.  The current
implementation choosing $1/(\gradrsq\Vfc)$ is consistent with the
paper's printed equation but inherits the paper's dimensional shorthand
(Assumption A9 above).

\section{Dimensional check of the rigorous form}
\label{sec:dimcheck}

With $[\Dped]=[f]=\mathrm{m^{2}/s}$, $[\Sreact]=\mathrm{m^{3}/s}$,
$[\Vfc]=[\Vcx]=\mathrm{m/s}$, $[n_{0}]=\mathrm{m^{-3}}$, $[L]=\mathrm{m}$,
$[\dneneg]=\mathrm{m^{-4}}$, $[\fFC]=[\fCX]=1$, $[\gradrsq]=1$,
$[\nFCzero]=\mathrm{m^{-3}}$, $[\mathcal{E}]=1$:
\begin{align}
  [A]   &\;=\; \frac{\mathrm{m^{-3}}\cdot\mathrm{m}\cdot\tfrac{\mathrm{m^{3}}}{\mathrm{s}}}
                   {\tfrac{\mathrm{m}}{\mathrm{s}}}
               \;=\; 1, \\
  [B]   &\;=\; \mathrm{m^{2}}\cdot\frac{\tfrac{\mathrm{m^{3}}}{\mathrm{s}}\cdot\mathrm{m^{-4}}}
                   {\tfrac{\mathrm{m}}{\mathrm{s}}}
               \;=\; 1, \\
  [\Ecx]&\;=\; \mathrm{m^{2}}\cdot\frac{\tfrac{\mathrm{m^{3}}}{\mathrm{s}}\cdot\mathrm{m^{-3}}}
                   {\tfrac{\mathrm{m^{2}}}{\mathrm{s}}}
               \;=\; 1, \\
  [K]   &\;=\; 1.
\end{align}
All four coefficients are dimensionless. \hfill$\checkmark$

\section{Signs of the terms in Eq.~(15)}
\label{sec:signs}

The sign structure mirrors Eqs.~(1)--(4), just with substitutions:
\begin{itemize}
  \item The overall $-\ne\Si[\cdot]$ prefactor is the \emph{electron sink
        due to ionization}, inherited from Eq.~(1).  Ionization always
        removes a neutral and adds an electron, so
        $\nabla\cdot\Gamma_{e}=+\ne\Si(\fsa{\nfc}+\fsa{\ncx})$; moving
        the minus sign inherent to $\nabla\cdot(-\Dped\nabla\ne)$ to the
        RHS gives the $-$ in front.
  \item Inside the bracket, the \emph{first term} $-\gradrsq\Dped(\dnexnow-\dneneg)/(\Vcx\fCX)$
        is the contribution of $\fsa{\ncx}$ from~\eqref{eq:12}.  Its
        sign enters as follows:  both $\dnexnow$ and $\dneneg$ are
        \emph{negative} in the pedestal (density decreases outward), and
        $\dnexnow \le \dneneg < 0$ (the pedestal gradient is steeper
        than the core), so $(\dnexnow-\dneneg)\le 0$.  Divided by
        $\Vcx\fCX>0$ and flipped by the overall $-$, this term is
        non-negative, which matches $\ncx\ge 0$.
  \item The \emph{second term} inside the bracket,
        $\Ccx\,\nFCzero\,\mathcal{E}$, is the residual (post-CX-subtraction)
        contribution of $\fsa{\nfc}$.  Normally $\Ccx>0$ because
        $\Vfc\fFC/(\Vcx\fCX)<1$ at the separatrix in realistic plasmas
        (FC atoms are faster than thermal $\Vcx$ in the pedestal foot),
        so this term is positive.  It represents the net ionization
        source from FC atoms, already reduced by the fraction of the FC
        flux that was converted to CX.
  \item The \emph{exponent} in $\mathcal{E}$ is \emph{negative} for $x<0$
        because $\int_{0}^{x}=-\int_{x}^{0}$ and the integrand
        $\ne(\Si+\Scx)/(\fFC\Vfc)>0$.  So $\mathcal{E}(x)<1$ for $x<0$,
        and $\mathcal{E}\to 0$ deep in the plasma.  This implements the
        physical fact that FC atoms are exponentially attenuated going
        inward.
\end{itemize}
Putting these together, for $x<0$:
$\ne\Si\{\text{bracket}\}\ge 0$, so
$\tfrac{d}{dx}[\gradrsq\Dped\ne']\le 0$.  Integrating from $-\infty$ to
$0$ then gives a net positive electron inflow at $x=0$, exactly the
``neutral penetration balance'' expected on physical grounds.

\section{Summary of assumptions (beyond Eq.~(16))}
\label{sec:assumptions}

Every assumption carried through the Eq.~(16) derivation is inherited
unchanged, except that \textbf{A6 (no CX)} is \emph{relaxed} in Eq.~(15).
We list them for completeness, and add one that is specific to Eq.~(15):

\begin{itemize}
  \item[\textbf{A1}] Locally mono-energetic neutrals; $\Vfc$ genuinely
        constant, $\Vcx=\Vcx(x)$ locally thermal.
  \item[\textbf{A2}] Slowly varying geometry and form factors ($\gradrsq$,
        $\oint R^{2}d\theta$, $\fFC$, $\fCX$ $\simeq$ const in $d/dx$).
  \item[\textbf{A2b}] Slowly varying $\Vcx$; drop $(d\Vcx/dx)\fsa{\ncx}$.
  \item[\textbf{A3}] Slab-like neutral closure: in going
        Eq.~(5)$\to$(8) and (6)$\to$(9) the explicit $\gradrsq$ factor on
        the neutral LHS is dropped (kept only in the electron transport
        operator).
  \item[\textbf{A4}] $\fFC=1$: poloidal symmetry of FC neutrals
        (\texttt{solver.py} line~520).
  \item[\textbf{A4b}] $\fCX=1$: poloidal symmetry of CX neutrals (same).
  \item[\textbf{A6 (relaxed)}] $\fsa{\ncx}$ retained.  Its value is
        determined algebraically by~\eqref{eq:12}.
  \item[\textbf{A7}] Neumann BC at $\xin$: $\ne'(\xin)=\dneneg$,
        $\fsa{\nfc}(\xin)\approx 0$, $\fsa{\ncx}(\xin)\approx 0$.
  \item[\textbf{A8}] Quasineutrality with $Z_{i}=1$.
  \item[\textbf{A9 (new)}] When writing Eq.~(15) in the compact paper
        form~\eqref{eq:15-paper}, the combination $\gradrsq/(\Vcx\fCX)$
        is absorbed into the ``$\Dped$'' symbol in the first term ---
        a dimensional shorthand, analogous to A5 for Eq.~(16).  The
        rigorous form~\eqref{eq:15-derived} does \emph{not} make this
        identification.
\end{itemize}

\paragraph{Iteration structure.}  The code solves~\eqref{eq:nondim-ODE}
iteratively (\texttt{solver.py} line~913 onward): the exponential
kernel~$\mathcal{E}(x)$ depends on $\ne(x)$ itself through the integrand
$\ne(\Si+\Scx)/(\fFC\Vfc)$, so Eq.~(15) is a non-local integro-differential
equation.  Starting from the no-CX solution (\texttt{first\_step},
Eq.~(16)), $\mathcal{E}$ is built from the previous $\ne$ and treated as
a known function for the next BVP solve.  The iteration converges when
the electron density profile stops changing between steps.

\end{document}
